\documentclass[12pt,a4paper]{article}

\usepackage[
  paperwidth=14.8cm,
  paperheight=21cm,
  margin=2cm
]{geometry}
\usepackage[T1]{fontenc}
\usepackage[utf8]{inputenc}
\usepackage{lmodern}
\usepackage{enumitem}
\usepackage{romannum}
\newcommand{\croman}[1]
{\MakeUppercase{\romannumeral #1}}

% ---------- Meta data (edit these) ----------
\newcommand{\ThesisTitle}{Structure-function relationship of the molecular chaperone DNAJB6b at atomistic resolution} % line break with %
\newcommand{\ThesisAuthor}{Vasista Adupa}

% ---------- Layout tweaks ----------
\setlength{\parindent}{0pt}
\setlength{\parskip}{0pt}

% Right–aligned chapter label under each proposition
\newcommand{\chapterref}[1]{%
 \par\hfill\textit{Chapter #1}%
}

\begin{document}
\thispagestyle{empty}

\begin{center}
 Propositions accompanying the PhD thesis\\[1.5ex]
 \textbf{\ThesisTitle}\\[1.5ex]
  by \textbf{\ThesisAuthor}
\end{center}

\vspace{3ex}

\begin{enumerate}[label=\arabic*., leftmargin=*, itemsep=1.8ex, topsep=1.0ex]

  \item DNAJB6b samples three conformational states: closed, open, and extended. The closed and open states dominate, while the extended state is rare.
\chapterref{2}

\item The autoinhibited state of DNAJB6b remains stable across conformational changes, indicating that release likely requires interaction with other protein.
\chapterref{2}

\item Efficient suppression of polyglutamine aggregation by DNAJB6b relies on a balance between its holdase activity and its interaction with Hsp70.
\chapterref{3}

\item Helix \croman{5} contains both hydrophobic and charged residues that regulate the stability and release of DNAJB6b autoinhibition, respectively.
Hydrophobic residues, particularly phenylalanine's (F), stabilize the autoinhibited state through interactions with helices \croman{2} and \croman{3}, whereas acidic residues, such as aspartic acid (D), interact with lysine residues in the C-terminal domain (CTD) to promote release from autoinhibition.
\chapterref{3}

\item Self-association of DNAJB6b partially counteracts its chaperoning function in suppressing polyglutamine aggregation and is mediated by the substrate-binding S/T domain, which is also responsible for the holdase activity of the molecular chaperone.
\chapterref{3,4}

\item Coarse-grained simulations show that DNAJB6b acts primarily during the early stages of polyQ aggregation and has little effect on preformed fibrils.
\chapterref{5}

\end{enumerate}

\end{document}
